The definitions of Section \ref{sec.orderDrivenMarkets} describe a book at one instant.
This section is about where those instants come from,
and about why a participant computes them rather than receiving them.

\subsection{What an exchange disseminates}

A matching engine produces one thing: a stream of events.
It publishes that stream in two forms.

The first is the \emph{message stream}:
one message per submission, cancellation, deletion and execution,
carrying the identity of the order it concerns.
It is the tape of the orders themselves,
and the trading epoch of Definition \ref{def.orderQueue} is exactly what it transmits.

The second is the \emph{aggregated book}:
the prices and the sizes resting at each of the first $L$ levels of each side,
republished as they change.
It is the configuration
$(\bestAskPrice_t, \bestBidPrice_t, \lbrace(\nthBestAskSize[i]_t, \nthBestBidSize[i]_t)\rbrace)$
of Section \ref{sec.orderDrivenMarkets}, truncated at depth $L$.\footnote{
Both are shipped as separate products on every major venue.
Eurex publishes the message stream as the Enhanced Order Book Interface,
and the aggregated book twice over:
as the Enhanced Market Data Interface, un-netted,
and as the Market Data Interface, whose updates are netted at regular intervals
\cite{DBG24eob, DBG24emd}.
On CME's MDP 3.0 the same event appears in a market-by-order and in a market-by-price book
\cite{CMEmdp, CME26mbo, CME26mbl},
and Nasdaq publishes the message stream as TotalView-ITCH \cite{NAS23itc, NAS26tot}.
Venues and vendors call the first of the two market by order, order by order or level three,
and the second market by price, depth of book or level two
\cite{B2Bmdp, ONXmdp, EPAmdp, DBNmdp, MDLitc}.
The nomenclature is not standardised.}

The two are not independent products.
The queues $\askOrderQueue_t$ and $\bidOrderQueue_t$ are sets of orders,
and the sizes $\nthBestAskSize[i]_t$ and $\nthBestBidSize[i]_t$ are sums over them,
so the second stream is the image of the first under those sums.
The exchange computes that image by applying the matching rules of
Section \ref{sec.orderDrivenMarkets} to its own message stream,
which is to say that it runs the fold of the next subsection before publishing the result.

A sum does not record its summands:
the aggregate says how much rests at a price,
never how that size is split between orders,
nor which of them the matching algorithm reaches first.
That split is what decides whether a given resting order is executed,
and Section \ref{sec.orderDrivenMarkets} sets it aside.

\begin{remark}\label{remark.whyReconstruct}
 The aggregated book retains the quantity this chapter is about:
 how size has accumulated across the price levels of the two sides.
 That cross-section is what the queue imbalance \eqref{eq.queueImbalance} and the order flow
 imbalance of Section \ref{sec.priceFormation} are computed from,
 and Section \ref{sec.priceFormation} is the claim that the movement of the price emerges
 from its statistics.
\end{remark}

\begin{remark}\label{remark.derivedAndDelayed}
 Notice that the aggregated book is derived, and that the derivation sits in the
 publication path:
 a message is on the wire once the engine has written it,
 whereas a level update exists only once the event has been netted against the level's
 total.
 It is also truncated at depth $L$, and where the feed is netted its updates are sent at
 regular intervals, so that intermediate configurations are absent rather than delayed.
 The aggregation is left out of the message stream by design,
 price levels and their aggregation being derivable by the recipient from the individual
 orders, and the stream being published
 ``in order to send public market data as fast as possible''
 \cite[Section 3.1]{DBG24eob}.
 A participant who folds the stream herself therefore holds the aggregated book sooner,
 and to a depth the feed does not publish.
\end{remark}

\begin{remark}\label{remark.aggregatedAsCheck}
 The two feeds therefore have different jobs.
 A fold run on the message stream is the participant's own computation:
 it is available sooner and to any depth,
 and a dropped message or a mishandled event leaves it wrong
 in a way that no statistic computed from it reveals.
 The aggregated book is the exchange's own fold of the same stream,
 published later and truncated,
 and comparing the two level by level locates such an error where nothing else does.
 It is a check on the reconstruction rather than a substitute for it.
\end{remark}

A second tape exists in United States equities for a regulatory rather than a technical
reason.
Rule 603(c) of Regulation NMS obliges anyone who displays quotations,
where a trading decision can be implemented,
to display consolidated data alongside them \cite{SEC05reg}.
Consolidation is performed by a party other than the venue on which the quotation arose,
and costs a latency that a participant reading each venue's own feed does not pay.

\subsection{The book as a fold}

Writing $\mathcal{B}$ for a configuration and $m_k$ for the $k$-th message of the stream,
\begin{equation}\label{eq.bookFold}
 \mathcal{B}_{k} = F(\mathcal{B}_{k-1}, m_k),
 \qquad \mathcal{B}_0 = \varnothing ,
\end{equation}
so that the book is the accumulator of a fold and holds one state, the current one.
The sequence $(\mathcal{B}_k)_k$ is not held by the book;
it is produced by whatever drives the fold, at whatever resolution is wanted.

The transition $F$ is Proposition \ref{prop.lobUpdate}.
What makes it a single function rather than a case analysis is
Proposition \ref{prop.decompositionOfLimitOrder}:
every incoming order is processed as its market part followed by its resting part,
so there is no separate branch for a marketable order,
and an implementation that writes one has two code paths where the mathematics has one.

A message is a submission, carrying an order $(t,q,p,d)$,
or a cancellation, withdrawing a size $q$ from the queue at $p$ on the side $d$;
a deletion is the cancellation of all of it.
An execution is not a third kind:
by Remark \ref{remark.passiveDirection} it is the record of the market part of somebody
else's submission, read on the passive side.
Algorithm \ref{algo.bookFold} is $F$.

\begin{algorithm}[h!]
 \caption{{The fold of the message stream into the aggregated book}}
 \label{algo.bookFold}
 \begin{algorithmic}[5]
  \REQUIRE $\mathcal{B}_{k-1}$ as two maps from price to size, and the message $m_k$
  \IF{$m_k$ is a cancellation of $q$ at $p$ on the side $d$}
  \STATE subtract $q$ from the size at $p$ on the side $d$, and remove $p$ if it reaches zero
  \ELSE
  \STATE let $(t,q,p,d)$ be the submitted order and set
         $\marketOrderSize$ as in Proposition \ref{prop.decompositionOfLimitOrder}
  \STATE set $r := \marketOrderSize$
  \WHILE{$r > 0$}
  \STATE let $\pi$ be the best price-eligible price on the side $-d$
  \STATE match $r \wedge \nthBestAskSize[1]$ shares at $\pi$ if $d=+1$, and
         $r \wedge \nthBestBidSize[1]$ if $d=-1$
  \STATE subtract them at $\pi$, remove $\pi$ if it reaches zero,
         and decrease $r$ by the same amount
  \ENDWHILE
  \STATE add $q - \marketOrderSize$ to the size at $p$ on the side $d$
  \ENDIF
  \RETURN $\mathcal{B}_{k}$
 \end{algorithmic}
\end{algorithm}

The loop is the reduction of Proposition \ref{prop.lobUpdate} to its elementary pieces,
and it terminates because each pass empties one price.
We implement it more than once,
and compare the implementations on the time and the memory they take.

\begin{remark}\label{remark.gridAgainstOccupied}
 The index $i$ of $\nthBestAskSize[i]$ counts levels of the price grid,
 whereas an aggregated feed reports the first $L$ \emph{occupied} prices of each side.
 The two agree only when the occupied prices are contiguous on the grid,
 so the queue imbalance \eqref{eq.queueImbalance} is recoverable from a truncated aggregated
 book only for those $n$ within which no gap falls.
\end{remark}

\subsection{LOBSTER}

LOBSTER \cite{HP11lob} reconstructs the limit order book of a Nasdaq-traded stock from
Nasdaq's historical TotalView-ITCH files and publishes the result at
\url{https://lobsterdata.com}.
It distributes, for each ticker and each trading day,
a \emph{message} file and an \emph{orderbook} file:
the historical, file-based form of the same pair of feeds.
The message file carries the time, the event type, an order identifier, a size, a price and a
direction;
the orderbook file carries $4L$ columns, the price and size of each side repeated outward
from the touch, padded with sentinels where a side has fewer than $L$ occupied prices.

Six event types occur:
the submission of a limit order,
its partial cancellation,
its total deletion,
the execution of a visible limit order,
the execution of a hidden one,
and a trading halt.

\begin{remark}\label{remark.passiveDirection}
 Notice that there is no event for an aggressive order.
 A trade is recorded as the execution of the \emph{resting} order it matched,
 so the direction carried by an execution names the passive side
 and the aggressor's direction is its opposite.
 In the decomposition of Proposition \ref{prop.decompositionOfLimitOrder}
 the message stream records the second component of the pair explicitly,
 and the first only through its effect on the opposite side.
 Signing a trade is therefore exact and free here,
 and inverted by anyone who reads the column as the direction of the trade.
\end{remark}

The two files are aligned by position alone:
row $k$ of the orderbook file is the configuration after the event on row $k$ of the message
file, and the orderbook file carries no timestamp and no key of its own.
Since the first file determines the second through \eqref{eq.bookFold},
the pair is the check of Remark \ref{remark.aggregatedAsCheck} in recorded form:
an implementation of the fold is compared against the exchange's own, state by state.
Section \ref{sec.lobsterEmpirics} works from these files,
and \cite{LPV22sho} is the same question asked of them at scale.
